\paragraph{Formal-group heights.}
Artin and Mazur attach a formal group to the top coherent line;
Stienstra expresses its logarithm by integral coefficient numerators.
For a smooth degree-$d$ hypersurface $X\subset\PP^{d-1}_{\Fp}$,
$d\ge3$, write $\Phi_X$ for its degree-$(d-2)$ Artin--Mazur formal
group and $h(X)$ for its height, allowing $h(X)=\infty$.
It is a smooth one-dimensional formal group, including the plane cubic
case through its formal Picard group. If $\widetilde F$ is any
homogeneous integral lift, its normalized logarithm is
\[
 L(T)=\sum_{m\ge1}b_m\frac{T^m}{m},\qquad
 b_m=[(x_1\cdots x_d)^{m-1}]\widetilde F^{m-1}.
\]
These are \citet[II, Corollaries~2.11--2.12 and~4.2,
pp.~106--110]{fg-artin-mazur} and
\citet[Theorem~1, p.~907]{fg-stienstra}. In particular the first height
is detected by $b_p\bmod p$, not by reduction of the possibly
nonintegral coefficient $b_p/p$.

\begin{theorem}[Fixed height predicates]\label{fg:thm:complexity-intro}
Fix a prime $p$. Inputs are division-free arithmetic circuits over $\Fp$
promised to compute homogeneous degree-$d$ polynomials in $d\ge3$
variables defining smooth geometrically integral projective
hypersurfaces, with dimension allowed to grow. The following statements
hold under deterministic polynomial-time many-one reductions.
\begin{enumerate}
\item Deciding $h(X)=1$ is $\ModP$-complete, and
      $h(X_\Phi)=1$ in Theorem~\ref{thm:main} precisely when
      $S\not\equiv0\pmod p$.
\item For every fixed integer $r\ge2$, deciding $h(X)\ge r$ is
      $\ModP$-complete; the predicate includes infinite height.
\item For every fixed finite $h\ge1$, deciding $h(X)=h$ is
      $\ModP$-complete. Hardness in these statements persists after
      imposing any fixed lower bound on the dimension.
\end{enumerate}
At fixed $p$, every fixed finite height predicate is polynomial time for
an explicit sparse equation. In fixed dimension the full height of a
circuit equation is polynomial-time computable. No polynomial bound in
a varying binary-encoded prime or threshold is asserted by the fixed
parameter statements.
\end{theorem}

The coefficient and reduction proofs are in
Sections~\ref{fg:sec:height}--\ref{fg:sec:space}.
Exact-height hardness uses an explicit prime-field existence theorem.

\begin{theorem}[Explicit realization of every finite height]
\label{fg:thm:seeds-intro}
For every prime $p$ and every finite integer $h\ge1$, there is an
explicit smooth geometrically integral hypersurface over $\Fp$ of
degree $p^h+p+1$ in $p^h+p+1$ homogeneous variables whose
one-dimensional top Artin--Mazur formal group has exact height $h$.
Its dimension is $p^h+p-1$, and its canonical bundle is trivial.
Its equation can be constructed in time polynomial in its displayed
degree, with $p$ and $h$ as parameters; no polynomial bound in the
binary length of a varying $h$ is asserted.
\end{theorem}

The construction in Section~\ref{fg:sec:seeds} prescribes a unique
multinomial contribution to each relevant logarithm numerator and then
preserves those coefficients under smoothing. It gives equations over
the prime field. The all-height construction of
\citet[Theorem~4.18]{fg-kty-two} works over an algebraically closed
field and uses a smaller degree; its nonempty parameter loci alone
do not specify a prime-field point.

The next height has an explicit counting interpretation. For a
width-three CNF $\phi$, use two variables $w_{f,i,0},w_{f,i,1}$ for
each clause-variable incidence. If $I_f$ is the set of variables in
clause $f$ and $R_f\subseteq\{0,1\}^{I_f}$ is its set of satisfying
local assignments, define
\[
 P_f=\sum_{a\in R_f}\prod_{i\in I_f}w_{f,i,a_i}.
\]
Put
\[
 P=\prod_fP_f\prod_i
       \left(\prod_{f\ni i}w_{f,i,0}+\prod_{f\ni i}w_{f,i,1}\right),
 \qquad Z=\prod_{(f,i)}w_{f,i,0}w_{f,i,1}.
\]
An unused declared variable receives a tautological unary relation.
If there is no incidence, first add a forced variable. A local
configuration chooses one term from each displayed factor, so
coefficient multiplicities distinguish these choices.

\begin{theorem}[The second height counts cyclic orbits]\label{fg:thm:orbit-intro}
Fix $p$, and let $S$ be the number of satisfying assignments of $\phi$.
Set $A_2(\phi)=[Z^{p+1}]P^{p+1}$. There is a function
$C(\phi)\in\#\mathrm P$ satisfying the exact integer identity
\[
 A_2(\phi)=S^2+pC(\phi).
\]
It counts the nonfixed cyclic orbits of balanced ordered $(p+1)$-tuples
of local configurations: their product monomial is $Z^{p+1}$, and the
cyclic group rotates the first $p$ configurations while fixing the last.
For the hypersurface of Theorem~\ref{thm:main},
\[
 h(X_\phi)\ge3
 \iff S=0\pmod p\ \text{and }C(\phi)=0\pmod p
 \iff A_2(\phi)=0\pmod{p^2}.
\]
On $S=0\bmod p$, its second normalized height invariant is
$b_{p^2}/p=C(\phi)\bmod p$. This invariant can differ for two CNF
presentations of the same Boolean function, even with the same output
degree and ambient dimension.
\end{theorem}

\begin{theorem}[The full height in polynomial space]\label{fg:thm:space-intro}
Fix $p$ and the circuit promise of Theorem~\ref{fg:thm:complexity-intro}.
If $L$ is the complete binary input length, the finite height can be
output in binary, or infinity identified, with
$O_p(L^2\log(L+2))$ working bits. Every finite height is at most
$\binom{pd-1}{d-1}$. The algorithm halts on all syntactically valid
inputs, with correctness required on the promise. In particular,
infinite-height decision belongs to polynomial space.
\end{theorem}

The geometric finite operator is due to
\citet[Theorem~C]{fg-kty-one}, and its matrix formulation to
\citet[\S3.4]{fg-bgp}. The space bound follows by accessing its entries
from the original circuit and recursively computing one matrix-power
entry at a time. A matching hardness classification for infinite height
is not established here.

\begin{theorem}[Integral derivatives of a height-$h$ lift]\label{fg:thm:derivative-intro}
Let $R=W(k)$ for a perfect characteristic-$p$ field, let $\mathcal F$ be
a one-dimensional commutative formal-group law over $R$, and let
$D=(L'(T))^{-1}d/dT$ be its normalized integral invariant derivation.
If the reduction has height at least $h\ge1$, put
\[
 c_h=\frac1{p-1}-\frac1{p^h-1},\qquad
 E_h(a)=\sum_{\substack{j\ge1\\h\nmid j}}\lfloor a/p^j\rfloor.
\]
For every $f\in R[[T]]$ and $a\ge0$,
\[
 D^af\in p^{\max\{0,E_h(a),\lceil c_ha-(h-1)\rceil\}}R[[T]].
\]
At exact height $h$, for $j=hr+s$, $0\le s<h$,
$v_p(D^{p^j}(T^{p^s})(0))=E_h(p^j)$. The slope $c_h$ and uniform
affine error $h-1$ are optimal. At infinite height the exact uniform
exponent is $v_p(a!)$.
For a Stienstra lift of height at least $r$, with $q=p^r$,
\[
 h(X)\ge r+1\iff b_q\in p^rR
      \iff D^qR[[T]]\subseteq q!R[[T]].
\]
\end{theorem}

Thus $S=0\bmod p$ implies the height-two derivative law, with slope
$p/(p^2-1)$ and affine error one, for every integral Stienstra lift over
$\Z_p$. In particular $\overline D^{p}=0$. The next height is detected
by the additional divisibility of $D^{p^2}T(0)$: its valuation is $p$
at height two and at least $p+1$ at height at least three. The generic
proofs are given in Section~\ref{fg:sec:derivatives}; directional
statements for a specified Honda Cartier dual require its additional
integral model data.

